\chapter{Tests et qualité logicielle}

\section{Stratégie de tests}

La stratégie de tests de \textbf{My Smart Budget} suit la pyramide de tests avec une approche pragmatique visant 80\% de couverture de code. L'accent est mis sur les fonctionnalités critiques : gestion des transactions financières, calcul des budgets, et sécurité des données.

\subsection{Architecture de tests multiniveaux}

\begin{tcolorbox}[colback=blue!5!white,colframe=blue!60!black,title=\textbf{Pyramide de tests My Smart Budget}]
\begin{verbatim}
                    +===================================+
                    |           TESTS E2E (4%)          |
                    |         (Cypress - 28 tests)      |
                    |    Parcours critiques complets    |
                    |    Temps: ~3 min / Nuit           |
                    +===================================+
                              |
                              |
                    +===================================+
                    |     TESTS D'INTEGRATION (13%)     |
                    |      (Jest/Supertest - 87 tests)  |
                    |    API endpoints, PostgreSQL      |
                    |       Temps: ~30s / Push          |
                    +===================================+
                              |
                              |
                    +===================================+
                    |      TESTS UNITAIRES (83%)        |
                    |        (Jest - 569 tests)         |
                    |    Services, Utils, Components    |
                    |     Temps: ~10s / Save            |
                    +===================================+
\end{verbatim}
\end{tcolorbox}

\subsection{Tests unitaires - Couverture détaillée}

L'exemple ci-dessous montre un test unitaire typique du service de calcul budgétaire. Il vérifie que les alertes sont déclenchées aux bons seuils et que les cas limites (budget à zéro, dépassement) sont correctement gérés :

\begin{exemple}
\textbf{Test unitaire du service de calcul budgétaire :}
\begin{lstlisting}[language=JavaScript]
// tests/unit/services/budget-calculator.test.js
const BudgetCalculatorService = require('../../../src/services/budget-calculator');
const Transaction = require('../../../src/models/Transaction');
const Budget = require('../../../src/models/Budget');

describe('BudgetCalculatorService', () => {
  let budgetCalculator;
  let mockTransactionModel;
  let mockBudgetModel;

  beforeEach(() => {
    // Mocks des accès base de données (pg-promise)
    mockTransactionModel = {
      query: jest.fn(),  // pg-promise db.any()
      queryOne: jest.fn() // pg-promise db.one()
    };
    mockBudgetModel = {
      query: jest.fn(),
      queryOne: jest.fn()
    };
    
    budgetCalculator = new BudgetCalculatorService(
      mockTransactionModel,
      mockBudgetModel
    );
  });

  describe('calculateRemainingBudget', () => {
    it('should calculate remaining budget correctly', async () => {
      // Arrange
      const budgetId = 1;
      const userId = 42;
      const budget = {
        id: budgetId,
        user_id: userId,
        monthlyLimit: 500,
        categoryId: 'cat789',
        startDate: new Date('2024-01-01')
      };
      
      const transactions = [
        { amount: 50, type: 'expense', date: new Date('2024-01-05') },
        { amount: 75, type: 'expense', date: new Date('2024-01-10') },
        { amount: 25, type: 'expense', date: new Date('2024-01-15') }
      ];

// ... (autres assertions de test omises pour concision)
\end{lstlisting}
\end{exemple}

Ces tests vérifient que le service de calcul se comporte correctement dans tous les cas, y compris les situations limites. Ils s'exécutent en moins de 10 secondes à chaque sauvegarde.

\subsection{Tests d'intégration - API et Base de données}

Les tests d'intégration valident que les couches communiquent bien entre elles. L'exemple suivant teste les endpoints de transactions de bout en bout, avec une vraie base de données de test :

\begin{exemple}
\textbf{Tests d'intégration des endpoints critiques :}
\begin{lstlisting}[language=JavaScript]
// tests/integration/api/transactions.test.js
const request = require('supertest');
const { db } = require('../../../src/db');
const app = require('../../../src/app');
const bcrypt = require('bcryptjs');
const jwt = require('jsonwebtoken');

describe('Transactions API Integration Tests', () => {
  let authToken;
  let testUser;
  let testCategory;

  beforeAll(async () => {
    // Création utilisateur de test (PostgreSQL)
    testUser = await db.one(
      `INSERT INTO users(email, password_hash, first_name, last_name)
       VALUES($1,$2,$3,$4) RETURNING *`,
      ['test@mysmartbudget.com', await bcrypt.hash('Test123!@#', 12),
       'Test', 'User']
    );
    
    // Génération token JWT
    authToken = jwt.sign(
      { userId: testUser.id, role: 'user' },
      process.env.JWT_SECRET,
      { expiresIn: '1h' }
    );
    
    // Création catégorie de test
    testCategory = await db.one(
      `INSERT INTO categories(name, icon, color) VALUES($1,$2,$3) RETURNING *`,
      ['Alimentation', 'shopping-cart', '#4CAF50']
    );
  });

  afterAll(async () => {
    await db.none('DELETE FROM transactions');
    await db.none('DELETE FROM categories');
    await db.none('DELETE FROM users');
    await db.$pool.end();
  });

  describe('POST /api/transactions', () => {
    it('should create a transaction with valid data', async () => {
      const transactionData = {
        amount: 45.50,
        type: 'expense',
        category_id: testCategory.id,
        description: 'Courses Carrefour',
        date: new Date().toISOString()

// ... (tests de filtrage, sécurité et isolation utilisateur omis)
\end{lstlisting}
\end{exemple}

\begin{exemple}
\textbf{Tests d'intégration  Rapports analytiques (MongoDB/Mongoose) :}
\begin{lstlisting}[language=JavaScript]
// tests/integration/api/reports.test.js
// Les rapports mensuels et logs d'audit sont stockes dans MongoDB
const mongoose = require('mongoose');
const MonthlyReport = require('../../../src/models/MonthlyReport');
const AuditLog = require('../../../src/models/AuditLog');

describe('Monthly Reports - MongoDB Integration', () => {
  beforeAll(async () => {
    await mongoose.connect(process.env.MONGO_URI_TEST);
  });

  afterAll(async () => {
    await MonthlyReport.deleteMany({});
    await AuditLog.deleteMany({});
    await mongoose.connection.close();
  });

  it('should persist monthly report with category breakdown', async () => {
    const report = await MonthlyReport.create({
      userId: 42, // FK vers PostgreSQL users.id
      period: { month: 1, year: 2025 },
      metrics: {
        totalExpenses: 1800,
        savingsRate: 28,
        categoryBreakdown: [
          { categoryName: 'Alimentation', amount: 450, percentage: 25 },
          { categoryName: 'Transport', amount: 180, percentage: 10 }
        ]
      }
    });

    expect(report._id).toBeDefined();
    expect(report.metrics.categoryBreakdown).toHaveLength(2);
    expect(report.metrics.savingsRate).toBe(28);
  });

  it('should log audit events for sensitive operations', async () => {
    await AuditLog.create({
      userId: 42,
      action: 'LOGIN_SUCCESS',
      metadata: { ip: '127.0.0.1', userAgent: 'Jest/Test' },
      timestamp: new Date()
    });

    const log = await AuditLog.findOne({ userId: 42, action: 'LOGIN_SUCCESS' });
    expect(log).toBeTruthy();
    expect(log.metadata.ip).toBe('127.0.0.1');
  });
});
\end{lstlisting}
\end{exemple}

\begin{tcolorbox}[colback=blue!5!white,colframe=blue!60!black,title=\textbf{Stratégie de tests Polyglot (PostgreSQL + MongoDB)}]
L'architecture dual-database implique deux stratégies de tests distinctes :
\begin{itemize}
    \item \textbf{pg-promise (PostgreSQL)} : tests des données critiques (transactions, budgets, utilisateurs)  contraintes ACID, intégrité référentielle, requêtes paramétrées.
    \item \textbf{Mongoose (MongoDB)} : tests des données non-structurées (rapports mensuels, logs d'audit)  schémas flexibles, agrégations analytiques, TTL index.
\end{itemize}
Les deux suites de tests s'exécutent en parallèle dans le pipeline CI/CD, chacune avec sa propre base de test isolée.
\end{tcolorbox}

\subsection{Tests End-to-End - Parcours utilisateur complets}

\begin{exemple}
\textbf{Test E2E Cypress - Parcours création de budget :}
\begin{lstlisting}[language=JavaScript]
// cypress/e2e/budget-creation.cy.js
describe('Budget Creation Flow', () => {
  beforeEach(() => {
    // Reset database et login
    cy.task('db:seed');
    cy.login('test@mysmartbudget.com', 'Test123!@#');
    cy.visit('/budgets');
  });

  it('should create a new budget envelope successfully', () => {
    // Vérifier la page des budgets
    cy.contains('h1', 'Mes Enveloppes Budgétaires').should('be.visible');
    
    // Cliquer sur le bouton de création
    cy.get('[data-cy=create-budget-btn]').click();
    
    // Remplir le formulaire
    cy.get('[data-cy=budget-name]').type('Alimentation');
    cy.get('[data-cy=budget-amount]').type('400');
    cy.get('[data-cy=budget-category]').select('Courses');
    cy.get('[data-cy=budget-color]').click();
    cy.get('[data-cy=color-green]').click();
    
    // Soumettre
    cy.get('[data-cy=submit-budget]').click();
    
    // Vérifier la création
    cy.contains('.budget-card', 'Alimentation').within(() => {
      cy.contains('400€').should('be.visible');
      cy.contains('0% utilisé').should('be.visible');
      cy.get('.progress-bar').should('have.css', 'width', '0px');
    });
    
    // Vérifier la notification
    cy.get('.notification-success')
      .should('contain', 'Enveloppe créée avec succès')
      .and('be.visible');
  });

  it('should update budget status when adding transactions', () => {
    // Créer un budget
    cy.createBudget('Loisirs', 200);
    
    // Ajouter une transaction
    cy.visit('/transactions');
    cy.get('[data-cy=add-transaction]').click();
    cy.get('[data-cy=amount]').type('50');
    cy.get('[data-cy=category]').select('Loisirs');
    cy.get('[data-cy=description]').type('Cinéma');
    cy.get('[data-cy=submit-transaction]').click();

// ... (tests alerte 80% et dépassement 100% omis)
\end{lstlisting}
\end{exemple}

\subsection{Métriques de couverture de tests}

\begin{tcolorbox}[colback=green!5!white,colframe=green!60!black,title=\textbf{Couverture de code actuelle (Décembre 2025)}]
\begin{center}
\begin{tabular}{|l|c|c|c|c|c|}
\hline
\textbf{Module} & \textbf{Statements} & \textbf{Branches} & \textbf{Functions} & \textbf{Lines} & \textbf{Tests} \\
\hline
\textbf{Services} & 91.2\% & 87.5\% & 94.1\% & 91.8\% & 142 \\
\hline
\textbf{Controllers} & 82.7\% & 79.3\% & 85.2\% & 83.1\% & 98 \\
\hline
\textbf{Models} & 78.4\% & 72.1\% & 80.0\% & 78.9\% & 67 \\
\hline
\textbf{Middleware} & 95.3\% & 92.8\% & 96.7\% & 95.5\% & 54 \\
\hline
\textbf{Utils} & 98.1\% & 96.4\% & 100\% & 98.2\% & 89 \\
\hline
\textbf{Components Next.js} & 76.8\% & 71.2\% & 78.3\% & 77.1\% & 234 \\
\hline
\hline
\textbf{TOTAL} & \textbf{85.3\%} & \textbf{81.2\%} & \textbf{87.4\%} & \textbf{85.7\%} & \textbf{684} \\
\hline
\end{tabular}
\end{center}

\textbf{Fichiers avec couverture critique (< 50\%) nécessitant attention :}
\begin{itemize}
    \item \texttt{src/utils/legacy-import.js} : 32\% (code legacy à refactorer)
    \item \texttt{src/controllers/admin.controller.js} : 48\% (tests en cours)
    \item \texttt{src/components/Charts/ComplexChart.jsx} : 41\% (difficile à tester)
\end{itemize}
\end{tcolorbox}

\section{Tests de performance}

\subsection{Stratégie de tests de charge}

Les tests de performance de My Smart Budget utilisent k6 pour simuler des charges réalistes correspondant aux pics d'utilisation mensuels (début de mois lors des imports bancaires).

\begin{exemple}
\textbf{Scénario de test de charge k6 - Import mensuel :}
\begin{lstlisting}[language=JavaScript]
// k6/scenarios/monthly-import-spike.js
import http from 'k6/http';
import { check, sleep, group } from 'k6';
import { SharedArray } from 'k6/data';
import { Rate, Trend } from 'k6/metrics';

// Métriques personnalisées
const loginErrors = new Rate('login_errors');
const importDuration = new Trend('import_duration');
const dashboardLoad = new Trend('dashboard_load_time');

// Données de test partagées
const testUsers = new SharedArray('users', function() {
  return JSON.parse(open('./test-users.json'));
});

export const options = {
  scenarios: {
    // Simulation pic mensuel d'imports
    monthly_spike: {
      executor: 'ramping-arrival-rate',
      startRate: 10,
      timeUnit: '1s',
      preAllocatedVUs: 10,
      maxVUs: 50,
      stages: [
        { duration: '2m', target: 5  },  // Montée progressive
        { duration: '3m', target: 25 }, // Pic d'activité
        { duration: '5m', target: 25 }, // Maintien charge élevée
        { duration: '2m', target: 5  },  // Redescente
        { duration: '2m', target: 2  },   // Retour normal
      ],
    }
  },
  thresholds: {
    http_req_duration: ['p(95)<800', 'p(99)<2000'], // 95% < 800ms
    http_req_failed: ['rate<0.05'],                  // < 5% erreurs
    login_errors: ['rate<0.02'],                     // < 2% échecs login
    import_duration: ['p(95)<5000'],                 // Import < 5s
    dashboard_load_time: ['p(95)<2000'],             // Dashboard < 2s

// ... (fonction de test principale avec groupes login/dashboard/budgets)
\end{lstlisting}
\end{exemple}

\subsection{Résultats des tests de performance}

\begin{tcolorbox}[colback=yellow!5!white,colframe=yellow!60!black,title=\textbf{Rapport de performance k6 - Simulation pic mensuel (1er janvier 2026)}]
\textit{(Sortie k6 -- résultats disponibles dans les artéfacts CI/CD)}

\textbf{Analyse des résultats :}
\begin{itemize}
    \item ✅ \textbf{Objectif P95 < 800ms :} Atteint (782ms)
    \item ✅ \textbf{Taux d'erreur < 5\% :} Atteint (1.58\%)
    \item ✅ \textbf{Dashboard < 2s :} Atteint (1.98s en P95)
    \item ⚠️ \textbf{P99 élevé :} 1842ms (optimisation nécessaire pour cas extrêmes)
\end{itemize}

En pratique, cela signifie que l'application répond en moins d'une seconde dans 95\,\% des cas, même lors des pics de charge du début de mois. Seul le percentile 99 dépasse 1,8 secondes, ce qui représente des cas très rares à optimiser en priorité.
\end{tcolorbox}

\subsection{Monitoring de performance en production}

\begin{tcolorbox}[colback=blue!5!white,colframe=blue!60!black,title=\textbf{Métriques de performance production (Décembre 2025)}]
\begin{center}
\begin{tabular}{|l|c|c|c|c|}
\hline
\textbf{Endpoint} & \textbf{P50 (ms)} & \textbf{P95 (ms)} & \textbf{P99 (ms)} & \textbf{Req/jour} \\
\hline
POST /auth/login & 82 & 156 & 298 & 168 \\
GET /transactions & 124 & 342 & 567 & 840 \\
POST /transactions & 98 & 218 & 412 & 280 \\
GET /dashboard & 256 & 687 & 1,234 & 560 \\
GET /budgets & 142 & 389 & 623 & 336 \\
\hline
\end{tabular}
\end{center}

\textbf{Optimisations appliquées suite aux tests :}
\begin{itemize}
    \item ✅ Requêtes SQL optimisées avec index composites (gain 40\% sur P95)
    \item ✅ Index PostgreSQL optimisés sur user\_id + date (gain 60\% sur queries)
    \item ✅ Pagination côté serveur (limit 50 transactions)
    \item ✅ Compression gzip activée (réduction 70\% payload)
    \item ✅ CDN pour assets statiques (gain 200ms sur First Paint)
\end{itemize}
\end{tcolorbox}

\section{Qualité du code avec SonarQube}

\subsection{Configuration et intégration CI/CD}

\begin{exemple}
\textbf{Configuration SonarQube pour My Smart Budget :}
\begin{lstlisting}
# sonar-project.properties
sonar.projectKey=my-smart-budget
sonar.projectName=My Smart Budget
sonar.projectVersion=1.0.0
sonar.organization=mysmartbudget

# Sources
sonar.sources=src
sonar.exclusions=**/*.test.js,**/*.spec.js,**/node_modules/**,**/coverage/**
sonar.tests=tests

# Couverture
sonar.javascript.lcov.reportPaths=coverage/lcov.info
sonar.testExecutionReportPaths=test-results/jest-results.xml

# Qualité
sonar.javascript.globals=React,process,__dirname,Buffer
sonar.javascript.environments=browser,node,jest

# Seuils de qualité (Quality Gate)
sonar.qualitygate.wait=true
sonar.qualitygate.timeout=300

# Règles personnalisées
sonar.issue.ignore.multicriteria=e1,e2
sonar.issue.ignore.multicriteria.e1.ruleKey=javascript:S3776
sonar.issue.ignore.multicriteria.e1.resourceKey=**/legacy/*.js
sonar.issue.ignore.multicriteria.e2.ruleKey=Web:BoldAndItalicTagsCheck
sonar.issue.ignore.multicriteria.e2.resourceKey=**/*.jsx
\end{lstlisting}
\end{exemple}

\subsection{Métriques de qualité actuelles}

\begin{tcolorbox}[colback=green!5!white,colframe=green!60!black,title=\textbf{Dashboard SonarQube - My Smart Budget (15 décembre 2025)}]
\begin{center}
\begin{figure}[H]
    \centering
    \safefig{chapters/assets/images/sonarqube-badge.png}{0.70\textwidth}{sonarqube-badge.png}
    \caption{Badge SonarQube -- résumé qualité du code (My Smart Budget)}
    \label{fig:sonarqube-badge}
\end{figure}
\end{center}

\begin{center}
\begin{tabular}{|l|c|c|c|}
\hline
\textbf{Métrique} & \textbf{Valeur} & \textbf{Objectif} & \textbf{Statut} \\
\hline
\textbf{Bugs} & 0 & 0 & ✅ Passed \\
\textbf{Vulnerabilities} & 0 & 0 & ✅ Passed \\
\textbf{Security Hotspots} & 3 reviewed & 100\% reviewed & ✅ Passed \\
\textbf{Code Smells} & 47 & < 100 & ✅ Passed \\
\textbf{Coverage} & 85.3\% & > 80\% & ✅ Passed \\
\textbf{Duplications} & 2.1\% & < 3\% & ✅ Passed \\
\textbf{Maintainability} & A & A & ✅ Passed \\
\textbf{Reliability} & A & A & ✅ Passed \\
\textbf{Security} & A & A & ✅ Passed \\
\hline
\end{tabular}
\end{center}

\textbf{Détail des Code Smells par catégorie :}
\begin{itemize}
    \item \textbf{Complexité cognitive :} 12 (fonctions trop complexes)
    \item \textbf{Code dupliqué :} 8 (à refactorer en utils)
    \item \textbf{Conventions naming :} 15 (variables mal nommées)
    \item \textbf{Dead code :} 6 (code non utilisé)
    \item \textbf{Console.log oubliés :} 6 (à retirer)
\end{itemize}
\end{tcolorbox}

\begin{exemple}
\textbf{Exemple de refactoring suite à SonarQube :}
\begin{lstlisting}[language=JavaScript]
// AVANT : Complexité cognitive élevée (score: 21)
const calculateBudgetStatus = (budget, transactions) => {
  let total = 0;
  let status = 'healthy';
  
  for (let i = 0; i < transactions.length; i++) {
    if (transactions[i].categoryId === budget.categoryId) {
      if (transactions[i].type === 'expense') {
        if (transactions[i].date >= budget.startDate && 
            transactions[i].date <= budget.endDate) {
          total += transactions[i].amount;
          
          if (total > budget.limit) {
            status = 'critical';
          } else if (total > budget.limit * 0.8) {
            status = 'warning';
          } else if (total > budget.limit * 0.5) {
            status = 'normal';
          }
        }
      }
    }
  }
  
  return { total, status, remaining: budget.limit - total };
};

// APRÈS : Complexité réduite (score: 8)
const calculateBudgetStatus = (budget, transactions) => {
  const relevantTransactions = filterRelevantTransactions(
    transactions, 
    budget
  );
  
  const totalSpent = calculateTotalSpent(relevantTransactions);
  const status = determineBudgetStatus(totalSpent, budget.limit);
  
  return {
    total: totalSpent,
    status,
    remaining: budget.limit - totalSpent,
    percentage: (totalSpent / budget.limit) * 100
  };
};

const filterRelevantTransactions = (transactions, budget) => {
  return transactions.filter(t => 
    t.categoryId === budget.categoryId &&
    t.type === 'expense' &&
    isWithinDateRange(t.date, budget.startDate, budget.endDate)
  );
};

const calculateTotalSpent = (transactions) => {
  return transactions.reduce((sum, t) => sum + t.amount, 0);
};

const determineBudgetStatus = (spent, limit) => {
  const percentage = (spent / limit) * 100;
  
  if (percentage > 100) return 'critical';
  if (percentage > 80) return 'warning';
  if (percentage > 50) return 'normal';
  return 'healthy';
};

const isWithinDateRange = (date, start, end) => {
  return date >= start && date <= end;
};
\end{lstlisting}
\end{exemple}

\subsection{Métriques Lighthouse}

\begin{tcolorbox}[colback=blue!5!white,colframe=blue!60!black,title=\textbf{Scores Lighthouse - My Smart Budget (Mobile 4G)}]
\begin{center}
\begin{tabular}{|l|c|c|c|}
\hline
\textbf{Catégorie} & \textbf{Score} & \textbf{Objectif} & \textbf{Détails} \\
\hline
\textbf{Performance} & 92/100 & > 90 & ✅ FCP: 1.2s, LCP: 2.1s, TTI: 2.8s \\
\textbf{Accessibilité} & 95/100 & > 90 & ✅ ARIA labels, contraste OK \\
\textbf{Best Practices} & 93/100 & > 85 & ✅ HTTPS, pas de console.log \\
\textbf{SEO} & 91/100 & > 85 & ✅ Meta tags, sitemap \\
\textbf{PWA} & 86/100 & > 80 & ✅ Manifest, offline basic \\
\hline
\end{tabular}
\end{center}

\textbf{Opportunités d'amélioration identifiées :}
\begin{itemize}
    \item ⚠️ Réduire le JavaScript inutilisé (économie potentielle : 142KB)
    \item ⚠️ Optimiser les images (format WebP, économie : 89KB)
    \item ⚠️ Précharger les fonts critiques (gain : 200ms)
    \item ⚠️ Implémenter le cache service worker complet
\end{itemize}
\end{tcolorbox}

\section{Synthèse et plan d'amélioration}

\begin{tcolorbox}[colback=green!5!white,colframe=green!60!black,title=\textbf{Points forts de la stratégie qualité}]
\begin{itemize}
    \item ✅ \textbf{Couverture de tests élevée :} 85.3\% (objectif 80\% dépassé)
    \item ✅ \textbf{0 bugs critiques} en production depuis 3 mois
    \item ✅ \textbf{Pipeline CI/CD robuste :} 94.3\% de succès
    \item ✅ \textbf{Performance optimisée :} P95 < 800ms sur tous les endpoints critiques
    \item ✅ \textbf{Qualité A} sur toutes les métriques SonarQube
    \item ✅ \textbf{Lighthouse > 90} sur Performance et Accessibilité
    \item ✅ \textbf{Tests automatisés :} 684 tests exécutés à chaque commit
\end{itemize}
\end{tcolorbox}

\begin{tcolorbox}[colback=yellow!5!white,colframe=yellow!60!black,title=\textbf{Plan d'amélioration Q1 2026}]
\textbf{Objectifs prioritaires :}
\begin{enumerate}
    \item \textbf{Augmenter la couverture à 90\%}  Mon objectif personnel est de concentrer l'effort sur les composants Next.js (actuellement 76,8\%) et d'ajouter des tests de mutation (Stryker.js).
    
    \item \textbf{Réduire le temps du pipeline à moins de 10 minutes}  Je prévois de paralléliser les tests E2E et d'activer le cache Docker pour ne pas freiner mon rythme de développement.
    
    \item \textbf{Améliorer la performance P99}  Objectif P99 < 1500ms (actuellement 1842ms) via l'optimisation des requêtes PostgreSQL (index, query plans) et la mise en cache HTTP sur les endpoints lourds.
    
    \item \textbf{Zero Code Smell}  Refactoring des 47 code smells restants identifiés par SonarQube.
\end{enumerate}
\end{tcolorbox}

